\section{Atmosphere model}\label{sec:atmosphere}
% Phase 1 (23 Nov-27 Dec 2026).

\subsection{Choice of standard atmosphere}
% ISA vs. U.S. Standard Atmosphere 1976 \cite{ussa1976}: why one is used as the source
% and in which altitude range the two are identical. Why a 1-D, altitude-only model
% (no latitude/longitude/season dependence) is the usual choice for ascent simulation.
\subsection{Geometric and geopotential altitude}
% What geopotential altitude is, why the standard is defined in it, the conversion
% formula, and the size of the difference at typical apogees.
\subsection{Layer structure and governing equations}
% Hydrostatic equation + ideal gas; gradient vs. isothermal layers; layer table.
% Full derivation in \cref{app:derivations}.
\subsection{Range of validity}
% Altitude range implemented and what happens outside it.
\subsection{Validation}
% Figure: model vs. tabulated values; max relative error. Full table in \cref{app:validation}.
% TODO (learning note): the reference values are read from the printed USSA 1976 tables,
% not recomputed from the defining equations, so the check is independent of the code
% (a wrong constant would otherwise sit in both the model and the test and go unnoticed).
